%% F03 \dash{} Logarytmy i skale logarytmiczne
%%
%% Zalążek. Poniższy opis jest kontraktem tego programu: co ma uzasadnić i co
%% ma dać czytelnikowi. Napisanie programu oznacza usunięcie bloku
%% \programstub{} i zastąpienie go programem. Jeśli po przerobieniu materiału
%% opis okaże się błędny, popraw go i odnotuj sprzeczność w CLAUDE.md w sekcji
%% *Resolved questions*.

\program{Logarytmy i skale logarytmiczne}
\label{prog:F03}

\programstub{%
\textit{[Opis do przetłumaczenia \dash{} patrz wydanie angielskie.]}

Argument: a logarithm turns multiplication into addition, and that is the
only reason anybody uses one. Payoff: probabilities of a 2,000-token
sequence multiply to something no float can hold, so we sum log-
probabilities instead; perplexity is \code{exp} of a mean log; a loss curve
is plotted on a log axis because the interesting part is the ratio, not the
difference. This is the most load-bearing Foundation program in the book.

\medskip
\textbf{Planowana długość:} 45 ramek.
}
